\documentclass[10pt, aspectratio=169]{beamer}
% Preámbulo modular para presentaciones Beamer (Modelación Estadística)
% Diseño y paleta institucional del Tecnológico de Monterrey

\documentclass[aspectratio=169,xcolor={dvipsnames,table}]{beamer}

\usepackage[utf8]{inputenc}
\usepackage[spanish,mexico]{babel}
\usepackage{amsmath,amssymb,amsfonts,mathtools}
\usepackage{graphicx}
\usepackage{listings}
\usepackage{booktabs}
\usepackage{tikz}
\usetikzlibrary{matrix,arrows,decorations.pathmorphing,shapes.geometric,positioning}

% Paleta Institucional Tecnológico de Monterrey
\definecolor{TecAzul}{HTML}{0039A6}       % Azul TEC: color primario institucional
\definecolor{TecAzulOscuro}{HTML}{1A2E51} % Azul marino oscuro
\definecolor{TecRojo}{HTML}{EC2661}       % Rojo de acento (paleta miTec)
\definecolor{TecGris}{HTML}{646464}       % Gris neutro para comentarios y subtítulos
\definecolor{TecFondoBloque}{HTML}{F4F6F9}% Fondo suave para bloques

% Configuración de Tema Beamer
\usetheme{Boadilla}
\usecolortheme[named=TecAzulOscuro]{structure}
\setbeamercolor{alerted text}{fg=TecRojo}
\setbeamercolor{palette primary}{bg=TecAzulOscuro,fg=white}
\setbeamercolor{palette secondary}{bg=TecAzul,fg=white}
\setbeamercolor{palette tertiary}{bg=TecAzulOscuro!80!black,fg=white}
\setbeamercolor{title}{fg=TecAzulOscuro,bg=TecFondoBloque}
\setbeamercolor{frametitle}{fg=TecAzulOscuro,bg=TecFondoBloque}

% Bloques personalizados elegantes
\setbeamercolor{block title}{bg=TecAzulOscuro,fg=white}
\setbeamercolor{block body}{bg=TecFondoBloque,fg=black}
\setbeamercolor{block title alerted}{bg=TecRojo,fg=white}
\setbeamercolor{block body alerted}{bg=TecRojo!8,fg=black}
\setbeamercolor{block title example}{bg=TecAzul,fg=white}
\setbeamercolor{block body example}{bg=TecAzul!8,fg=black}

% Redondear bordes y añadir sombra sutil
\setbeamertemplate{blocks}[rounded][shadow=true]
\setbeamertemplate{navigation symbols}{}

% Pie de página limpio con número de diapositiva
\setbeamertemplate{footline}{
  \leavevmode%
  \hbox{%
  \begin{beamercolorbox}[wd=.7\paperwidth,ht=2.25ex,dp=1ex,left]{author in head/foot}%
    \hspace*{2ex}\insertshortauthor~~(\insertshortinstitute) \hfill \insertshorttitle
  \end{beamercolorbox}%
  \begin{beamercolorbox}[wd=.3\paperwidth,ht=2.25ex,dp=1ex,right]{date in head/foot}%
    \insertshortdate{}\hspace*{2em}
    \insertframenumber{} / \inserttotalframenumber\hspace*{2ex}
  \end{beamercolorbox}}%
  \vskip0pt%
}

% Configuración de Listings de Python para visualización pedagógica de código
\lstset{
    language=Python,
    basicstyle=\ttfamily\footnotesize,
    keywordstyle=\color{TecAzulOscuro}\bfseries,
    stringstyle=\color{TecRojo},
    commentstyle=\color{TecGris}\itshape,
    showstringspaces=false,
    breaklines=true,
    frame=lines,
    rulecolor=\color{TecAzulOscuro!30},
    backgroundcolor=\color{TecFondoBloque},
    aboveskip=0.5em,
    belowskip=0.5em,
    literate={á}{{\'a}}1 {é}{{\'e}}1 {í}{{\'i}}1 {ó}{{\'o}}1 {ú}{{\'u}}1
             {Á}{{\'A}}1 {É}{{\'E}}1 {Í}{{\'I}}1 {Ó}{{\'O}}1 {Ú}{{\'U}}1
             {ñ}{{\~n}}1 {Ñ}{{\~N}}1 {¿}{{?`}}1 {¡}{{!`}}1
}

% Metadatos institucionales por defecto
\title[Modelación Estadística]{Modelación Estadística para Ciencia de Datos e Ingeniería}
\author[J. Castillo Colmenares]{Juliho David Castillo Colmenares}
\institute[Tec de Monterrey]{Tecnológico de Monterrey \\ Escuela de Ingeniería y Ciencias}
\date{\today}

% Comandos y macros estadísticas compartidas para las presentaciones Beamer (Metropolis)

% Conjuntos numéricos básicos
\newcommand{\R}{\mathbb{R}}
\newcommand{\N}{\mathbb{N}}
\newcommand{\Q}{\mathbb{Q}}
\newcommand{\Z}{\mathbb{Z}}
\newcommand{\set}[1]{\left\{ #1 \right\}}
\newcommand{\abs}[1]{\left|#1\right|}
\newcommand{\norm}[1]{\left\| #1 \right\|}

% Comandos estadísticos y probabilísticos del curso
\newcommand{\Var}{\operatorname{Var}}
\newcommand{\cov}{\operatorname{Cov}}
\newcommand{\corr}{\rho}
\newcommand{\comb}[2]{\begin{pmatrix} #1 \\ #2 \end{pmatrix}}
\newcommand{\s}{\sigma}
\newcommand{\particion}{\mathfrak{P}}
\newcommand{\card}[1]{\#\left( #1 \right)}
\newcommand{\seq}[3]{\set{#1}_{#2}^{#3}}
\newcommand{\E}{\mathbb{E}}
\newcommand{\Prob}{\mathbb{P}}

% Entornos pedagógicos y cajas en español académico natural
\newenvironment{discusion}{%
  \begin{alertblock}{Pregunta de Discusión}%
}{%
  \end{alertblock}%
}

\newenvironment{reflexion}{%
  \begin{alertblock}{Reflexión para el Aula}%
}{%
  \end{alertblock}%
}

\newenvironment{paradoja}{%
  \begin{alertblock}{Paradoja de Exploración}%
}{%
  \end{alertblock}%
}

\newenvironment{motivacion}{%
  \begin{exampleblock}{Motivación: Ciencia de Datos en la Práctica}%
}{%
  \end{exampleblock}%
}

\newenvironment{retoclase}{%
  \begin{block}{Reto Rápido en Equipo}%
}{%
  \end{block}%
}


\title[Validación de Modelos]{Sección 09.08: Validación de Modelos}
\subtitle{División Entrenamiento/Prueba y Validación Cruzada de $k$ Iteraciones}
\author[J. Castillo Colmenares]{Juliho Castillo Colmenares}
\institute{Tecnológico de Monterrey}
\date{\vspace{-1.2cm}}

\begin{document}

% Slide 1: Portada
\begin{frame}[plain]
  \vspace{-0.8cm}
  \titlepage
\end{frame}

% Slide 2: Hoja de Ruta
\begin{frame}{Hoja de Ruta --- Capítulo 09: Regresiones Lineales y Múltiples}
  \scriptsize
  ¿Dónde nos encontramos en el desarrollo modular del capítulo?
  \begin{itemize}\itemsep=0.02em
    \item \textbf{09.01} Correlación como Premisa de la Regresión
    \item \textbf{09.02} Introducción a la Regresión Lineal
    \item \textbf{09.03} Matemáticas de la Regresión: Derivación MCO y $R^2$
    \item \textbf{09.04} Regresión Lineal sobre Datos Simulados
    \item \textbf{09.05} Coeficientes Óptimos, Pruebas $t$/$F$ y RSE
    \item \textbf{09.06} Implementación en Python con \texttt{statsmodels}
    \item \textbf{09.07} Regresión Múltiple, Ecuación Normal y Ridge/Lasso
    \item \textbf{09.08} Validación de Modelos y $k$-fold Cross-Validation \emph{(Hoy)}
    \item \textbf{09.09} Diagnóstico de Residuos y Multicolinealidad (VIF)
    \item \textbf{09.10} Regresión con \texttt{scikit-learn}
    \item \textbf{09.11} Variables Categóricas y Variables Muda
    \item \textbf{09.12} Transformaciones No Lineales y Regresión Polinomial
  \end{itemize}
  \pause
  \vspace{-0.05cm}
  \begin{block}{Objetivo de la Sesión}
    Determinar si un modelo ajustado \textbf{generaliza} a datos nuevos: implementar la división entrenamiento/prueba para detectar sobreajuste, construir la validación cruzada de $k$ iteraciones desde cero, y usarla para comparar modelos competidores de forma robusta.
  \end{block}
\end{frame}

% Slide 3: Motivación
\begin{frame}{¿Generaliza el Modelo, o Solo Memoriza?}
  \footnotesize
  \begin{motivacion}
    Un modelo predictivo necesita validarse observando su rendimiento en \textbf{datos distintos} a los usados para ajustarlo. Un modelo puede ajustarse extremadamente bien a los datos de entrenamiento (\emph{sobreajuste}, \emph{overfitting}) y aun así fallar completamente al predecir observaciones nuevas.
  \end{motivacion}

  \vspace{0.15cm}
  \pause
  Todos los diagnósticos de las Secciones 09.03-09.07 (incluyendo $R^2$, pruebas $t$/$F$) se calculan \textbf{dentro} de la muestra de ajuste: ninguno, por sí solo, garantiza buen desempeño fuera de ella.
\end{frame}

% Slide 4: División entrenamiento/prueba
\begin{frame}{División entre Datos de Entrenamiento y Prueba}
  \footnotesize
  \begin{block}{Procedimiento}
    Dividimos aleatoriamente el conjunto de datos (típicamente 80\%/20\%): ajustamos el modelo \textbf{solo} con los datos de entrenamiento, y evaluamos su error \textbf{solo} con los datos de prueba, nunca vistos durante el ajuste.
  \end{block}
  \pause

  \vspace{0.1cm}
  \begin{alertblock}{Señal de Sobreajuste}
    Si el error (RSE/RMSE) en el conjunto de prueba es sustancialmente mayor que en el de entrenamiento, el modelo memorizó peculiaridades de la muestra de ajuste en lugar de capturar la relación real subyacente.
  \end{alertblock}
\end{frame}

% Slide 5: k-fold cross-validation
\begin{frame}{Validación Cruzada de $k$ Iteraciones}
  \footnotesize
  \begin{block}{Procedimiento}
    Dividimos los datos en $k$ subconjuntos (\emph{folds}) de tamaño aproximadamente igual. En cada una de las $k$ iteraciones: entrenamos con $k-1$ folds y evaluamos en el fold restante. Promediamos las $k$ métricas obtenidas.
  \end{block}
  \pause

  \vspace{0.1cm}
  \begin{alertblock}{Ventajas sobre una Sola División}
    Menor varianza (cada partición específica influye menos en la estimación final); uso eficiente de datos (cada observación se usa exactamente una vez para prueba); detecta sobreajuste de forma más robusta que una única división.
  \end{alertblock}
\end{frame}

% Slide 6: Elección de k
\begin{frame}{Elección de $k$ y Validación Cruzada Estratificada}
  \footnotesize
  \begin{itemize}\itemsep=0.1cm
    \item $k=5$ o $k=10$: valores convencionales, buen balance sesgo-varianza. \pause
    \item $k$ pequeño: mayor sesgo, menor varianza, más rápido. \pause
    \item $k=n$ (\emph{leave-one-out}, LOOCV): casi insesgado, pero lento y de alta varianza. \pause
    \item \textbf{CV estratificada:} preserva la distribución de una variable de interés en cada fold, útil con datos desbalanceados.
  \end{itemize}
\end{frame}

% Slide 7: Lab Python (Bloque 1)
\begin{frame}[fragile]{Laboratorio en Python: Detección de Sobreajuste}
  \scriptsize
  Un modelo sobreparametrizado (16 predictores para $n=40$) ajusta bien el entrenamiento pero falla en prueba (\texttt{09.08\_model\_validation.py}):
  {\lstset{basicstyle=\fontsize{3.0pt}{3.7pt}\selectfont\ttfamily,aboveskip=0.05em,belowskip=0.05em}
  \lstinputlisting[firstline=36, lastline=62]{../../code/09_regresiones/09.08_model_validation.py}}
\end{frame}

% Slide 8: Lab Python (Bloque 2)
\begin{frame}[fragile]{Laboratorio en Python: $k$-Fold desde Cero}
  \scriptsize
  Implementación completa de validación cruzada de 5 iteraciones sin librerías externas:
  {\lstset{basicstyle=\fontsize{2.9pt}{3.5pt}\selectfont\ttfamily,aboveskip=0.05em,belowskip=0.05em}
  \lstinputlisting[firstline=65, lastline=85]{../../code/09_regresiones/09.08_model_validation.py}}
\end{frame}

% Slide 9: Lab Python (Bloque 3)
\begin{frame}[fragile]{Laboratorio en Python: Comparación de Modelos vía CV}
  \scriptsize
  Comparación robusta entre un modelo de un predictor y uno de dos predictores:
  {\lstset{basicstyle=\fontsize{3.2pt}{3.9pt}\selectfont\ttfamily,aboveskip=0.05em,belowskip=0.05em}
  \lstinputlisting[firstline=106, lastline=119]{../../code/09_regresiones/09.08_model_validation.py}}
\end{frame}

% Slide 10: Salida Terminal
\begin{frame}[fragile]{Laboratorio en Python: Salida en Terminal}
  \scriptsize
  Salida estándar generada al ejecutar el script del laboratorio en Python:
  \vspace{0.02cm}
  \begin{block}{Salida Estándar en Consola}
    \fontsize{3.4pt}{4.1pt}\selectfont
    \begin{verbatim}
=== Block 1: Overfitting Detection ===
n_train=32, n_test=8, p=16 (1 real + 15 pure noise)
RSE train=8.0786, RMSE test=12.4452 (diff=4.3666) -> overfitting signature

=== Block 2: k-Fold CV from Scratch ===
Folds: 0.9056, 0.8083, 0.8765, 0.8452, 0.8409
Mean R^2 = 0.8553 +/- 0.0332

=== Block 3: Comparing Models via CV ===
Model 1 (x1 only):  mean R^2 = 0.5727 +/- 0.0804
Model 2 (x1+x2):    mean R^2 = 0.8553 +/- 0.0332
    \end{verbatim}
  \end{block}
\end{frame}

% Slide 11: Interpretación
\begin{frame}{Interpretación: ¿Por Qué Importa la Validación?}
  \footnotesize
  \begin{itemize}\itemsep=0.12cm
    \item El Bloque 1 muestra el peligro de tener \textbf{más predictores que información real}: 15 columnas de ruido puro reducen artificialmente el error de entrenamiento, pero el error de prueba lo delata. \pause
    \item El Bloque 3 confirma, con evidencia \emph{fuera de muestra}, que el segundo predictor genuino mejora sustancialmente el modelo ($R^2$ promedio de $0.57$ a $0.86$) --- una conclusión más confiable que comparar $R^2$ dentro de la muestra. \pause
    \item La desviación estándar entre folds ($\pm0.033$ vs. $\pm0.080$) también revela que el modelo completo es más \textbf{estable}, no solo más preciso en promedio.
  \end{itemize}
\end{frame}

% Slide 12: Cierre
\begin{frame}{Cierre de Sección: Validación de Modelos}
  \begin{reflexion}
    Ningún diagnóstico calculado dentro de la muestra de ajuste --- ni $R^2$, ni las pruebas $t$/$F$ --- garantiza que un modelo generalice. La validación cruzada de $k$ iteraciones es el estándar robusto de la industria para estimar honestamente el desempeño de un modelo sobre datos que nunca ha visto.
  \end{reflexion}

  \vspace{0.2cm}
  \pause
  \begin{block}{Lecciones Clave}
    \begin{itemize}\itemsep=0.08cm
      \item \textbf{División entrenamiento/prueba:} rápida pero de alta varianza; útil para exploración inicial.
      \item \textbf{$k$-fold CV:} promedia sobre $k$ particiones, reduciendo la varianza de la estimación del error.
      \item \textbf{Comparación de modelos:} la CV es más confiable que comparar $R^2$ dentro de la muestra.
    \end{itemize}
  \end{block}
\end{frame}

% Slide 13: Perspectiva Modular
\begin{frame}{Perspectiva Modular --- El Próximo Reto}
  \scriptsize
  \begin{retoclase}
    Hemos validado el modelo \textbf{externamente} (generalización a datos nuevos). La Sección 09.09 lo audita \textbf{internamente}: verificaremos los supuestos clásicos sobre los errores (normalidad, homoscedasticidad, independencia) y cuantificaremos la multicolinealidad mediante el Factor de Inflación de Varianza --- dos formas de confiabilidad completamente distintas mas complementarias.
  \end{retoclase}

  \vspace{0.08cm}
  \pause
  \begin{alertblock}{Preparación Recomendada}
    \begin{itemize}\itemsep=0.06cm
      \item Reflexiona sobre por qué un modelo puede generalizar razonablemente bien (buena CV) y aun así violar los supuestos clásicos de regresión.
      \item Repasa la Matriz Sombrero de la Sección 09.07: su diagonal será clave para el diagnóstico de observaciones influyentes.
    \end{itemize}
  \end{alertblock}
\end{frame}

\end{document}
